\chapter{결과}
\label{ch:results}

\ifdraftmode
\begin{center}
\fcolorbox{red}{red!5}{%
\begin{minipage}{0.93\linewidth}
\small
\textcolor{red}{\textbf{본 장의 상태}}\par\smallskip
제 \ref{sec:kernel-results}절을 제외한 본 장의 모든 표는 \textbf{비어 있다}.
구현은 완성되었으나 실행되지 않았으므로(제 \ref{sec:impl-status}절) RQ1--RQ6의 어떤
수치도 존재하지 않는다. 미측정 항목은 \NM 로 표시하였다.\par\smallskip
표들은 손으로 채우도록 설계되지 않았다. \texttt{cvdyn2dgs experiments}가 산출하는
JSON을 \texttt{scripts/make\_tables.py}가 읽어 \texttt{paper/tables/}의 \texttt{.tex}
파일을 생성하며, 본 장은 그것을 \texttt{\textbackslash input}한다. 숫자를 원고에
직접 타이핑하는 경로는 존재하지 않는다.
\end{minipage}}
\end{center}
\medskip
\fi

\section{실행된 검증: 이산 논리 커널}
\label{sec:kernel-results}

본 절은 본 논문에서 유일하게 \emph{측정된} 결과이다. 제
\ref{sec:kernel-verification}절에서 기술한 커널 검증을 실행한 결과 41개 항목 전부가
통과하였다. PyTorch를 필요로 하지 않으므로 어떤 환경에서도 재현할 수 있다.

%% AUTO-GENERATED by scripts/make_tables.py -- DO NOT EDIT BY HAND.
%% Source: scripts/verify_kernels.py
%% Status: filled from measurements
\begin{table}[htbp]
\centering\small
\caption{실행된 커널 검증 결과. \texttt{python scripts/verify\_kernels.py}로 재현 가능.}
\label{tab:kernel-results}
\begin{tabular}{@{}lrr@{}}
\toprule
검증 묶음 & 항목 수 & 통과 \\
\midrule
큐브 $\to$ 6면체 분해            & 5  & 5 \\
marching tetrahedra 참조표      & 2  & 2 \\
(surfel, 타일) 페어 열거        & 3  & 3 \\
타일별 깊이 정렬과 절단          & 4  & 4 \\
점유 비트마스크 팩/언팩          & 3  & 3 \\
타원체 부호 거리 이분법          & 5  & 5 \\
삼선형 보간 가중치               & 3  & 3 \\
공액경사법                      & 2  & 2 \\
Godunov 상류 재초기화            & 2  & 2 \\
저장량 손익분기 대수             & 6  & 6 \\
수축 프로파일                   & 6  & 6 \\
\midrule
\textbf{합계}                   & \textbf{41} & \textbf{41} \\
\bottomrule
\end{tabular}
\end{table}


이 결과의 범위는 제한적이다. 통과는 이산 색인 논리가 건전함을 뜻하며, 파이프라인이
동작함을 뜻하지 않는다. 특히 래스터화, 자동미분, Chan--Vese 솔버, 전체 파이프라인은
전혀 포함되지 않는다.

\section{이론 수렴률}
\label{sec:results-rates}

제 \ref{ch:error}장의 수렴률 주장은 $\spacing$을 변화시키며 $\log$--$\log$ 기울기를
적합하여 확인한다.

%% AUTO-GENERATED by scripts/make_tables.py -- DO NOT EDIT BY HAND.
%% Source: results/theory_checks.json
%% Status: STUB - no measurements available
\begin{table}[htbp]
\centering\small
\caption{이론이 예측하는 수렴률과 측정된 경험적 지수.}
\label{tab:theory-rates}
\begin{tabular}{@{}lccl@{}}
\toprule
결과 & 예측 & 측정 & 비고 \\
\midrule
명제 \ref{prop:normal-error} $\sin\theta=O(\hmax^2)$
  & 2.0 & \NM & $\spacing$ 스윕 기울기 \\
보조정리 \ref{lem:quadratic} 이차 수축
  & 2.0 & \NM & 잔차 이력 적합 \\
명제 \ref{prop:interp-eikonal} 식 \eqref{eq:interp-eikonal}
  & 0.0 & \NM & 측정 대 예측 편차 \\
명제 \ref{prop:lowrank} Eckart--Young
  & 0.0 & \NM & 상대 편차 \\
\bottomrule
\end{tabular}
\end{table}


\pending{%
명제 \ref{prop:normal-error}이 예측하는 $\sin\theta = O(\hmax^2)$의 경험적 지수,
보조정리 \ref{lem:quadratic}이 예측하는 이차 수축의 경험적 차수, 명제
\ref{prop:planar-disk}의 $O(\kappa s^2)$ 비례상수, 명제 \ref{prop:interp-eikonal}의
측정 대 예측 eikonal 노름. 실행 명령: \texttt{cvdyn2dgs theory}.}

\section{RQ1: 웜스타트}
\label{sec:results-rq1}

%% AUTO-GENERATED by scripts/make_tables.py -- DO NOT EDIT BY HAND.
%% Source: results/research_questions.json
%% Status: STUB - no measurements available
\begin{table}[htbp]
\centering\small
\caption{RQ1: 웜스타트(식 \eqref{eq:warmstart-init})의 효과.}
\label{tab:rq1}
\begin{tabular}{@{}lcc@{}}
\toprule
 & 웜스타트 & 콜드스타트 \\
\midrule
총 반복 횟수            & \NM & \NM \\
프레임당 평균 반복      & \NM & \NM \\
총 시간 (ms)            & \NM & \NM \\
$K(e_0)$ 평균           & \NM & \NM \\
초기 오차 $e_0$         & \NM & \NM \\
평균 crop 비율          & \NM & \NM \\
\bottomrule
\end{tabular}
\end{table}


\pending{%
웜스타트와 콜드스타트의 총 반복 횟수 및 벽시계 시간, 참조 해에 대한 $K(e_0)$(식
\eqref{eq:iteration-bound}), 초기 오차 $e_0$. 정리 \ref{thm:warmstart}은 \emph{순서}만
예측하므로, 확인할 것은 $K(e^{\mathrm{warm}}_0)\le K(e^{\mathrm{cold}}_0)$이다.
벽시계 시간의 격차는 좁은 띠 crop 효과와 섞이므로 반복 횟수와 분리하여 해석해야 한다.}

\section{RQ2: 정사영 대 독립 재최적화}
\label{sec:results-rq2}

%% AUTO-GENERATED by scripts/make_tables.py -- DO NOT EDIT BY HAND.
%% Source: results/research_questions.json
%% Status: STUB - no measurements available
\begin{table}[htbp]
\centering\scriptsize
\setlength{\tabcolsep}{3.5pt}
\caption{RQ2: 법선 정사영 대 프레임별 독립 재최적화.}
\label{tab:rq2}
\begin{tabular}{@{}lrrrrrrrrrrr@{}}
\toprule
방법 & Dice & HD95 & $\Esurf$ & IoU & depth & normal & PSNR & SSIM & CR & 전처리 & p95 \\
     &      & (mm) & (mm)     &     & (mm)  & (deg)  & (dB) &      &    & (ms)   & (ms) \\
\midrule
CV-Dyn2DGS & \NM & \NM & \NM & \NM & \NM & \NM & \NM & \NM & \NM & \NM & \NM \\
Independent 2DGS & \NM & \NM & \NM & \NM & \NM & \NM & \NM & \NM & \NM & \NM & \NM \\
\bottomrule
\end{tabular}
\end{table}

\noindent 전처리 속도 배율: \NM$\times$.
프레임당 갱신 시간: \NM\,ms (제안) 대
\NM\,ms (재최적화).


\pending{%
식 \eqref{eq:cost-hypothesis}의 좌변 네 항($C_{\mathrm{CV}}$, $C_{\mathrm{project}}$,
$C_{\mathrm{orient}}$, $C_{\mathrm{res}}$)을 각각, 그리고 우변 $C_{\mathrm{full}}$을
측정하여 부등식이 성립하는지 확인한다. 품질 측면에서는 independent 2DGS를 상한으로
두고 비열등성 마진 안에 들어오는지 본다. 목표는 이기는 것이 아니라 마진 안에 있으면서
비용이 작은 것이다.}

\section{RQ3: 저장량}
\label{sec:results-rq3}

%% AUTO-GENERATED by scripts/make_tables.py -- DO NOT EDIT BY HAND.
%% Source: results/research_questions.json
%% Status: STUB - no measurements available
\begin{table}[htbp]
\centering\small
\caption{RQ3: 저장량과 압축비(식 \eqref{eq:s-full}--\eqref{eq:cr}). 세 가지 표면
저장 방식 모두를 보고한다.}
\label{tab:rq3}
\begin{tabular}{@{}lrrrr@{}}
\toprule
표면 저장 방식 & $\Sfull$ (MB) & $\Sours$ (MB) & $\CR$ & 손익분기 (B/프레임) \\
\midrule
좁은 띠 부호 거리 & \NM
  & \NM
  & \NM
  & \NM \\
binary mask & \NM
  & \NM
  & \NM
  & \NM \\
triangle mesh & \NM
  & \NM
  & \NM
  & \NM \\
\midrule
실제 파일 크기 (MB) & \multicolumn{4}{l}{\NM} \\
\bottomrule
\end{tabular}
\end{table}


\pending{%
$\Sfull$, $\Sours$, $\CR$을 세 가지 표면 저장 방식에 대해, 그리고 이론적 바이트 수와
실제 파일 크기를 나란히. 식 \eqref{eq:breakeven}의 손익분기 예산과 실제 프레임당 표면
크기의 비교가 핵심이다. $\CR<1$이면 제안 표현이 이 구성에서 이기지 못한다는 뜻이며,
그것도 유효한 결과이다.}

\section{RQ4: 재생 속도}
\label{sec:results-rq4}

%% AUTO-GENERATED by scripts/make_tables.py -- DO NOT EDIT BY HAND.
%% Source: results/research_questions.json
%% Status: STUB - no measurements available
\begin{table}[htbp]
\centering\small
\caption{RQ4: 재생 latency 분해(식 \eqref{eq:frame-budget}). 성공 기준은 p95 $<33.3$\,ms.}
\label{tab:rq4}
\begin{tabular}{@{}lrrrrrrrc@{}}
\toprule
해상도 & $N$ & surface & project & orient & raster & 평균 & p95 & 30 FPS \\
       &     & (ms)    & (ms)    & (ms)   & (ms)   & (ms) & (ms) &        \\
\midrule
256 & 10000 & \NM & \NM & \NM & \NM & \NM & \NM & \NM \\
256 & 20000 & \NM & \NM & \NM & \NM & \NM & \NM & \NM \\
512 & 10000 & \NM & \NM & \NM & \NM & \NM & \NM & \NM \\
512 & 20000 & \NM & \NM & \NM & \NM & \NM & \NM & \NM \\
\bottomrule
\end{tabular}
\end{table}


\pending{%
식 \eqref{eq:frame-budget}의 네 항을 분해하여, surfel 수와 viewer 해상도를 스윕하며
95번째 백분위 latency가 33.3\,ms 미만인지 확인한다. 또한 제 \ref{sec:seek-problem}절의
순차 대 정준 정사영 격차를 측정한다. 이 격차가 크면 임의 탐색을 식 \eqref{eq:store}의
저장량만으로 제공할 수 없다는 뜻이다.}

\section{RQ5: mesh와의 비교}
\label{sec:results-rq5}

%% AUTO-GENERATED by scripts/make_tables.py -- DO NOT EDIT BY HAND.
%% Source: results/research_questions.json
%% Status: STUB - no measurements available
\begin{table}[htbp]
\centering\small
\caption{RQ5: 동일 Chan--Vese 표면에서 2DGS와 텍스처 mesh의 지표별 비교. 단일 점수로
합치지 않고 지표별 승자를 보고한다.}
\label{tab:rq5}
\begin{tabular}{@{}lrrll@{}}
\toprule
지표 & 2DGS & mesh & 방향 & 승자 \\
\midrule
silhouette IoU & \NM & \NM & 높을수록 & \NM \\
boundary F-score & \NM & \NM & 높을수록 & \NM \\
depth RMSE (mm) & \NM & \NM & 낮을수록 & \NM \\
normal error (deg) & \NM & \NM & 낮을수록 & \NM \\
ROI PSNR (dB) & \NM & \NM & 높을수록 & \NM \\
ROI SSIM & \NM & \NM & 높을수록 & \NM \\
$\Eflicker$ & \NM & \NM & 낮을수록 & \NM \\
\bottomrule
\end{tabular}
\end{table}


\pending{%
동일 표면, 비슷한 primitive 예산에서 silhouette(문턱값 스윕), boundary F-score, depth
RMSE, normal angular error, ROI PSNR/SSIM, $\Eflicker$를 지표별로 비교한다. 단일 점수로
합치지 않고 지표별 승패를 보고한다. 2DGS가 어느 지표에서도 개선하지 못하면 ``이
목적에는 mesh가 더 적절하다''가 올바른 결론이다.}

\section{RQ6: 2D 원반 대 얇은 3D Gaussian}
\label{sec:results-rq6}

%% AUTO-GENERATED by scripts/make_tables.py -- DO NOT EDIT BY HAND.
%% Source: results/research_questions.json
%% Status: STUB - no measurements available
\begin{table}[htbp]
\centering\small
\caption{RQ6: 2D 원반과 얇은 3D Gaussian. $\rho$는 법선방향 두께 비.}
\label{tab:rq6}
\begin{tabular}{@{}lrrrr@{}}
\toprule
표현 & IoU & depth RMSE (mm) & normal (deg) & ROI PSNR (dB) \\
\midrule
2DGS (2D 원반) & \NM & \NM & \NM & \NM \\
\midrule
thin-3DGS rho=0.05 & \NM & \NM & \NM & \NM \\
thin-3DGS rho=0.1 & \NM & \NM & \NM & \NM \\
thin-3DGS rho=0.3 & \NM & \NM & \NM & \NM \\
\bottomrule
\end{tabular}
\end{table}


\pending{%
법선방향 두께 비 $\rho$를 스윕하며 depth RMSE와 normal angular error를 비교한다.
thin-3DGS의 깊이는 자락 전역에서 상수이므로 2DGS의 깊이 우위는 명제
\ref{prop:perspective}의 직접 검증이며 독립적 발견이 아니다. 더 얇은 타원체는
기하학적으로 원반에 가깝지만 조건수가 나빠지므로 그 절충이 스윕에서 드러난다.}

\section{기준선 종합}
\label{sec:results-baselines}

%% AUTO-GENERATED by scripts/make_tables.py -- DO NOT EDIT BY HAND.
%% Source: results/baselines.json
%% Status: STUB - no measurements available
\begin{table}[htbp]
\centering\scriptsize
\setlength{\tabcolsep}{3.5pt}
\caption{기준선 종합 비교 (표 \ref{tab:baselines}).}
\label{tab:baselines-results}
\begin{tabular}{@{}lrrrrrrrrrrr@{}}
\toprule
방법 & Dice & HD95 & $\Esurf$ & IoU & depth & normal & PSNR & SSIM & CR & 전처리 & p95 \\
     &      & (mm) & (mm)     &     & (mm)  & (deg)  & (dB) &      &    & (ms)   & (ms) \\
\midrule
cv-dyn2dgs & \NM & \NM & \NM & \NM & \NM & \NM & \NM & \NM & \NM & \NM & \NM \\
independent-2dgs & \NM & \NM & \NM & \NM & \NM & \NM & \NM & \NM & \NM & \NM & \NM \\
mesh-only & \NM & \NM & \NM & \NM & \NM & \NM & \NM & \NM & \NM & \NM & \NM \\
thin-3dgs & \NM & \NM & \NM & \NM & \NM & \NM & \NM & \NM & \NM & \NM & \NM \\
param-copy & \NM & \NM & \NM & \NM & \NM & \NM & \NM & \NM & \NM & \NM & \NM \\
closest-point & \NM & \NM & \NM & \NM & \NM & \NM & \NM & \NM & \NM & \NM & \NM \\
normal-no-residual & \NM & \NM & \NM & \NM & \NM & \NM & \NM & \NM & \NM & \NM & \NM \\
\bottomrule
\end{tabular}
\end{table}


\pending{%
표 \ref{tab:baselines}의 모든 기준선에 대한 주요 지표. 실행 명령:
\texttt{cvdyn2dgs experiments --size small}.}

\section{Ablation}
\label{sec:results-ablation}

%% AUTO-GENERATED by scripts/make_tables.py -- DO NOT EDIT BY HAND.
%% Source: results/ablations.json
%% Status: STUB - no measurements available
\begin{table}[htbp]
\centering\scriptsize
\setlength{\tabcolsep}{3.5pt}
\caption{Ablation (제 \ref{sec:ablation-design}절). 전체 설정에서 한 번에 하나의 기제만
토글한다.}
\label{tab:ablations}
\begin{tabular}{@{}lrrrrrr@{}}
\toprule
설정 & $\Esurf$ (mm) & IoU & depth (mm) & normal (deg) & PSNR (dB) & 전처리 (ms) \\
\midrule
full & \NM & \NM & \NM & \NM & \NM & \NM \\
no-warm-start & \NM & \NM & \NM & \NM & \NM & \NM \\
kp=1 & \NM & \NM & \NM & \NM & \NM & \NM \\
kp=2 & \NM & \NM & \NM & \NM & \NM & \NM \\
kp=3 & \NM & \NM & \NM & \NM & \NM & \NM \\
kp=5 & \NM & \NM & \NM & \NM & \NM & \NM \\
closest-point & \NM & \NM & \NM & \NM & \NM & \NM \\
isotropic & \NM & \NM & \NM & \NM & \NM & \NM \\
no-mask-loss & \NM & \NM & \NM & \NM & \NM & \NM \\
no-normal-loss & \NM & \NM & \NM & \NM & \NM & \NM \\
no-dist-loss & \NM & \NM & \NM & \NM & \NM & \NM \\
no-geometry-losses & \NM & \NM & \NM & \NM & \NM & \NM \\
no-residual & \NM & \NM & \NM & \NM & \NM & \NM \\
no-temporal-reg & \NM & \NM & \NM & \NM & \NM & \NM \\
full-rank-residual & \NM & \NM & \NM & \NM & \NM & \NM \\
lowrank-r4 & \NM & \NM & \NM & \NM & \NM & \NM \\
no-repulsion & \NM & \NM & \NM & \NM & \NM & \NM \\
no-densify & \NM & \NM & \NM & \NM & \NM & \NM \\
no-density-control & \NM & \NM & \NM & \NM & \NM & \NM \\
no-curvature-scale & \NM & \NM & \NM & \NM & \NM & \NM \\
\bottomrule
\end{tabular}
\end{table}


\pending{%
제 \ref{sec:ablation-design}절의 각 항목. 특히 간격 인지 여부의 ablation은 첫 번째
기여의 유일한 직접 시험이므로, 등방 격자와 비등방 격자에서의 격차를 함께 제시해야
한다.}

\section{개발 단계}
\label{sec:results-stages}

%% AUTO-GENERATED by scripts/make_tables.py -- DO NOT EDIT BY HAND.
%% Source: results/progressive.json
%% Status: STUB - no measurements available
\begin{table}[htbp]
\centering\scriptsize
\setlength{\tabcolsep}{3.5pt}
\caption{개발 단계별 결과 (표 \ref{tab:stages}).}
\label{tab:stages-results}
\begin{tabular}{@{}lrrrrrrrrrrr@{}}
\toprule
방법 & Dice & HD95 & $\Esurf$ & IoU & depth & normal & PSNR & SSIM & CR & 전처리 & p95 \\
     &      & (mm) & (mm)     &     & (mm)  & (deg)  & (dB) &      &    & (ms)   & (ms) \\
\midrule
v1-minimal & \NM & \NM & \NM & \NM & \NM & \NM & \NM & \NM & \NM & \NM & \NM \\
v2-geometry & \NM & \NM & \NM & \NM & \NM & \NM & \NM & \NM & \NM & \NM & \NM \\
v3-adaptive & \NM & \NM & \NM & \NM & \NM & \NM & \NM & \NM & \NM & \NM & \NM \\
\bottomrule
\end{tabular}
\end{table}


\pending{%
v1 $\to$ v2 $\to$ v3의 품질, 비용, 저장량 변화. 각 단계가 무엇을 얻고 무엇을
지불하는지를 함께 보아야 하며, 추가 계산만으로 산 향상은 그것으로 드러난다.}
